\documentclass[preview]{standalone}

\usepackage{amsmath}
\usepackage{amssymb}
\usepackage{stellar}
\usepackage{definitions}

\begin{document}

\id{field-extensions-exercises}
\genpage

\section{Exercises}

\begin{snippetexercise}{cubic-field-quotient-inverse-exercise}{A cubic field quotient}
    Let
    \[
        f=X^3-3X^2-X-1\in\rationalnumbers[X].
    \]
    Show that
    \[
        A=\rationalnumbers[X]/(f)
    \]
    is a \field, and find the inverse of
    \[
        (X^2+1)+(f)
    \]
    in \(A\).
\end{snippetexercise}

\begin{snippetsolution}{cubic-field-quotient-inverse-solution}{A cubic field quotient}
    A reducible cubic over a \field has a root. By the rational-root
    criterion, the only possible rational roots of \(f\) are \(1\) and
    \(-1\), but
    \[
        f(1)=-4,
        \qquad
        f(-1)=-4.
    \]
    Thus \(f\) is irreducible, so \((f)\) is a maximal \ideal and
    \(A\) is a field.

    Put \(g=X^2+1\). The Euclidean algorithm gives
    \[
        f=(X-3)g-2(X-1),
        \qquad
        g=(X-1)(X+1)+2.
    \]
    Hence
    \begin{align*}
        1
        &=
        \frac12g-\frac12(X-1)(X+1)\\
        &=
        \left(\frac12-\frac14(X-3)(X+1)\right)g
        +\frac14(X+1)f.
    \end{align*}
    Passing to the quotient yields
    \[
        \bigl((X^2+1)+(f)\bigr)^{-1}
        =
        \left(
            -\frac14X^2+\frac12X+\frac54
        \right)+(f).
    \]
\end{snippetsolution}

\end{document}
